\begin{frame}{Rappel: les 3 "niveaux" de l'écriture numérique (Bouchardon et Petit)}
    \begin{enumerate}
        \item La machine (binaire);
        \item Le programmeur (code);
        \item L'utilisateur (langage naturel).
    \end{enumerate}

\textrightarrow{} Sur le web, nous sommes tous écrivains.\\
\vspace{1em}
Le côté "graphique" de l'écriture numérique (écrire un mail, éditer un profil etc\dots) // l'écriture des données.
\end{frame}

\begin{frame}{La concurrence identitaire (Louise Merzeau)}
    \begin{center}
        {\small
        \textit{La personne qui sert de filtre à l’information n’est pas celle de nos CV, de nos appartenances et de nos papiers d’identité (même si, de plus en plus, elle en tient lieu). C’est une entité purement numérique qui se résume à la collection des traces laissées par nos connexions : requêtes, téléchargements, géolocalisation, achats, mais aussi contenus produits, copiés, repris, etc. Nous n’avons qu’une connaissance approximative de cette identité disséminée sur les réseaux.}}\\
        \textbf{Louise Merzeau, \textit{"Du signe à la trace, l'information sur mesure"}, Hermès no 53, 2009.}
    \end{center}
\end{frame}

\begin{frame}{Une écriture d'analphabètes (Cécile Portier)}
    \begin{center}
        {\small Si nous sommes les premiers producteurs des traces numériques qui nous concernent, nous en sommes très peu auteurs. Les données produites à notre insu sont quantitativement plus importantes que celles que nous pensons maîtriser, dans un discours, une rhétorique consciente et personnelle. De plus nous sommes des "auteurs" quasiment analphabètes : nous avons peu accès, nous ne savons que très peu lire les données qui nous définissent. Et pourtant, d'autres savent lire et exploiter ces données}\\
        \textbf{Cécile Portier, Traques traces}
    \end{center}
\end{frame}

\begin{frame}{Une écriture des plateformes : le sujet formaté}
    \textbf{Rappel plateformisation} = effet de standardisation et d'uniformisation de la construction identitaire.\\
    \vspace{1em}
Sur le plan graphique, les plateformes sont standardisées et souvent polluées par des éléments publicitaire contrairement au blogue.
\end{frame}